\documentclass[aspectratio=169,10pt,t]{beamer}
\usepackage{ac1}
\aulainfo{12}{Organização e Módulos}{Módulos em Python, reuso e estruturação de programas maiores}

\begin{document}
\capa

\begin{frame}{Roteiro}
  \begin{itemize}
    \item O que é um módulo? (Cap.~4.2)
    \item Formas de importar: \texttt{import}, \texttt{from}, \texttt{as}
    \item Criando seu próprio módulo
    \item \texttt{if \_\_name\_\_ == "\_\_main\_\_"}
    \item Estruturando um programa maior (4.6)
  \end{itemize}
\end{frame}

\begin{frame}{Módulos: bibliotecas de funções}
  \begin{columns}[T]
    \column{0.52\textwidth}
    \begin{itemize}
      \item Um \textbf{módulo} é um arquivo \texttt{.py} com funções e constantes prontas
      \item Python vem com a \emph{biblioteca padrão}: \texttt{math}, \texttt{random}, \texttt{statistics}, \texttt{datetime}, \texttt{os}\ldots
      \item Milhares de módulos de terceiros: \texttt{numpy}, \texttt{pandas}, \texttt{matplotlib}\ldots
    \end{itemize}
    \column{0.44\textwidth}
    \begin{exampleblock}{Filosofia}
      ``Não reinvente a roda'': antes de escrever, procure se já existe no manual do Python.
    \end{exampleblock}
  \end{columns}
\end{frame}

\begin{frame}[fragile]{Três formas de importar}
\begin{lstlisting}
import math                       # usa math.sin(...)
print(math.sin(math.pi / 2))      # 1.0

from math import radians, sin     # traz só estes nomes
print(sin(radians(30)))           # 0.49999999999999994

import math as m                  # apelido curto
print(m.sqrt(16))                 # 4.0
\end{lstlisting}
  \begin{itemize}
    \item \texttt{import math}: o prefixo \texttt{math.} deixa claro \emph{de onde} vem cada função
    \item \texttt{from \ldots import \ldots}: nomes diretos, mas risco de \emph{colisão} com seus nomes
    \item Evite \texttt{from math import *}
  \end{itemize}
\end{frame}

\begin{frame}[fragile]{Exemplo: tabela de senos e cossenos}
\begin{lstlisting}
from math import radians, sin, cos

print('{:>9}{:>9}{:>9}{:>9}'.format('Graus', 'Radianos',
                                    'Seno', 'Cosseno'))
for graus in range(0, 361, 30):
    rad = radians(graus)
    print('{:>9.2f}{:>9.2f}{:>9.2f}{:>9.2f}'.format(
        graus, rad, sin(rad), cos(rad)))
\end{lstlisting}
  \begin{itemize}
    \item Exercício 4.1: expandir a tabela para incluir a tangente (\texttt{tan})
    \item Exercício 4.4: pesquise no manual outras funções de \texttt{math}
  \end{itemize}
  \fonte{code/Cap4/prog.4.7\_tabela\_senos\_cossenosimport.py}
\end{frame}

\begin{frame}[fragile]{Criando seu próprio módulo}
  Arquivo \texttt{geometria.py}:
\begin{lstlisting}
import math

def area_circulo(raio):
    """Área do círculo."""
    return math.pi * raio ** 2

def distancia(x1, y1, x2, y2):
    """Distância entre dois pontos."""
    return math.hypot(x2 - x1, y2 - y1)
\end{lstlisting}
  Em outro arquivo, na mesma pasta:
\begin{lstlisting}
import geometria
print(geometria.area_circulo(5))
\end{lstlisting}
\end{frame}

\begin{frame}[fragile]{\texttt{\_\_name\_\_}: módulo ou programa?}
\begin{lstlisting}
# fim de geometria.py
def main():
    print("Teste:", area_circulo(1), distancia(0, 0, 3, 4))

if __name__ == "__main__":
    main()
\end{lstlisting}
  \begin{itemize}
    \item Executando \texttt{python geometria.py}: \texttt{\_\_name\_\_} vale \texttt{"\_\_main\_\_"} e o teste roda
    \item Com \texttt{import geometria}: \texttt{\_\_name\_\_} vale \texttt{"geometria"} e o teste \alert{não} roda
    \item Assim o mesmo arquivo serve como \textbf{programa} e como \textbf{biblioteca}
  \end{itemize}
\end{frame}

\begin{frame}{Estruturando programas maiores}
  \begin{columns}[T]
    \column{0.5\textwidth}
    \begin{block}{Camadas sugeridas}
      \begin{enumerate}\small
        \item \textbf{Entrada/saída}: lê do usuário/arquivos e imprime
        \item \textbf{Regras} (cálculos): funções puras
        \item \textbf{Dados}: listas, dicionários, arquivos
      \end{enumerate}
    \end{block}
    \column{0.47\textwidth}
    \begin{block}{Organização de arquivos}
      \ttfamily\footnotesize
      projeto/\\
      \hspace*{1em}principal.py\\
      \hspace*{1em}calculos.py\\
      \hspace*{1em}entrada\_saida.py\\
      \hspace*{1em}dados/
    \end{block}
  \end{columns}
  \vspace{1ex}
  Cada módulo tem \textbf{uma responsabilidade}. Isso será essencial no projeto final.
\end{frame}

\begin{frame}[fragile]{Exemplo: $\pi$ por Monte Carlo em módulos}
  \begin{columns}[T]
    \column{0.48\textwidth}
    \textbf{\texttt{montecarlo.py}}
\begin{lstlisting}[basicstyle=\ttfamily\scriptsize]
import random as r
import math as m

def gera_pontos():
    return r.random(), r.random()

def dentro(x, y):
    return m.hypot(x, y) < 1

def estima_pi(n):
    cont = 0
    for _ in range(n):
        x, y = gera_pontos()
        if dentro(x, y):
            cont += 1
    return 4 * cont / n
\end{lstlisting}
    \column{0.48\textwidth}
    \textbf{\texttt{principal.py}}
\begin{lstlisting}[basicstyle=\ttfamily\scriptsize]
import montecarlo

n = int(input("Pontos: "))
pi = montecarlo.estima_pi(n)
print("pi aprox.", pi)
\end{lstlisting}
    \vspace{1ex}
    \footnotesize Os testes e a entrada/saída ficam fora do módulo de cálculo.\\
    \fonte{code/Cap4/prog.4.23\_calcula\_pi\_main.py}
  \end{columns}
\end{frame}

\begin{frame}{Onde procurar ajuda}
  \begin{itemize}
    \item \texttt{help(math)} e \texttt{dir(math)} no interpretador
    \item Documentação oficial: \texttt{docs.python.org/pt-br/3/library/}
    \item Mensagem \texttt{ModuleNotFoundError}: o módulo não está instalado ou o arquivo não está na mesma pasta
    \item Instalar pacotes de terceiros: \texttt{pip install nome}
  \end{itemize}
  \begin{exampleblock}{Atividade}
    Crie um módulo \texttt{conversoes.py} com as funções \texttt{celsius\_para\_fahrenheit}, \texttt{fahrenheit\_para\_celsius} e \texttt{km\_para\_milhas}. Importe-o em outro arquivo e imprima uma tabela.
  \end{exampleblock}
\end{frame}

\exercicios{4}{%
  \begin{itemize}\small
    \item Módulos: \textbf{4.1} (tabela com tangente), \textbf{4.2} (raízes quadradas), \textbf{4.3} (quadrado e cubo), \textbf{4.4} (funções de \texttt{math})
  \end{itemize}}

\resumo{
  \item Módulo = arquivo \texttt{.py} reutilizável; \texttt{import} o traz para o seu programa
  \item \texttt{import x}, \texttt{from x import y}, \texttt{import x as z}
  \item Todo arquivo pode ser módulo; \texttt{if \_\_name\_\_ == "\_\_main\_\_"} protege o código de teste
  \item Programas grandes: módulos com responsabilidades claras
  \item Próxima aula: \texttt{math}, \texttt{random} e tratamento de erros
}

\end{document}
